\section*{Chapter \ref{ch:s1:asia-specific-equity-strategies} --- Asia-Specific Equity Strategies}

\begin{solution}{exo:s1:asia-specific-equity-strategies:1}
$\ln 1.5/\ln 1.1 = 4.25$, so five days under a 10\% limit, four of them locked; $\ln 1.5/\ln 1.2 = 2.22$, so three days under a 20\% limit, two locked.
\end{solution}

\begin{solution}{exo:s1:asia-specific-equity-strategies:2}
The backlog is $\ln 1.15 - \ln 1.10 = 0.044$ in log terms, so the fill probability is $e^{-0.044/0.02} = 0.108$: about one order in nine fills.
\end{solution}

\begin{solution}{exo:s1:asia-specific-equity-strategies:3}
$0.03 \times 0.6^4 = 0.0039$ remains, so $0.03 - 0.0039 = 0.0261$, 2.6 points in log terms, has reverted. The simulation's week-after figure, $-2.81\%$,
adds the fair value's own drift and noise.
\end{solution}

\begin{solution}{exo:s1:asia-specific-equity-strategies:4}
A pushed lock closes above the stock's value: sellers are glad to sell at the limit, so every queued order fills. The next open moves back toward the
value, which takes back part of the attention premium: 1.67\% against 5.90\% for locks with a backlog.
\end{solution}

\begin{solution}{exo:s1:asia-specific-equity-strategies:5}
A wider limit is reached only by larger moves, and the part of a large move that the limit holds back is larger: the backlog behind each lock grows with
the width. The attention premium is the same for all widths.
\end{solution}

\begin{solution}{exo:s1:asia-specific-equity-strategies:6}
$0.0129 \times 1{,}000 = 12.9$ a day. The filled trade earns $2.64\% - 0.15\% = 2.49\%$ on average, but only on the orders that fill, and those fill most
in the weakest locks.
\end{solution}

\begin{solution}{exo:s1:asia-specific-equity-strategies:7}
Limit-up closes fall to 0.79\% of stock-days; the next-open gain to 3.42\% for every close and 1.21\% after fills; the week's reversal to $-0.27\%$. The
backlog alone is worth about a point to a queued buyer; the rest of the trade, and nearly all the reversal, comes from the attention buyers.
\end{solution}

\begin{solution}{exo:s1:asia-specific-equity-strategies:8}
A close at the limit is a price at which almost nobody who wanted to buy could: the queue fills in the weakest locks and not in the strongest. Weighted
by fills the gain is 2.64\%, and the model's clean attention makes even that an upper bound. A Sharpe ratio needs fills from the queue, costs and the
capacity of the queue.
\end{solution}

\begin{solution}{pb:s1:asia-specific-equity-strategies:1}
\begin{enumerate}
\item A locked stock's further move in the same direction at the next open; prices accelerating toward a limit as they approach it.
\item They cut the observed variance, make returns positively autocorrelated and put point masses at the limits.
\item Selling shares the day they were bought; buying at the limit and selling the next day.
\item Fair returns scaled by 1.5; closes within the limit; attention buying of 3\% after limit-up closes, decaying by 0.6 a day; half the targets within
2\% of the limit pushed to lock; queue fills falling with the backlog.
\item Attention buying by individuals, a rise followed by reversal over the week, and 1.16\% a day to smart traders who buy on the limit day and sell the
next.
\item They buy on limit-up days and sell the next; their limit-day buying predicts a stronger long-run reversal; their pushing explains the \omterm{def:s1:asia-specific-equity-strategies:magnet}{magnet effect}.
\item Small accounts cannot predict returns and chase the day's moves; large accounts predict correctly and trade contrarian, but costs absorb it.
\item Northbound net purchases predict the returns of connected Shanghai stocks.
\item 11.24\%, 1.29\% and 0.10\% of stock-days.
\item 4.51\% and 2.64\%.
\item The queue fills most when the lock is weakest and least when the backlog is largest.
\item It moves days from the band 7.5\%--9.5\% to the limit: limit-up closes rise from 0.87\% to 1.29\% of stock-days.
\item \textbf{Named result.} At the next open after a limit-up close, a buyer at every close gains 2.98\%, 4.51\% and 6.58\% under 5\%, 10\% and 20\%
limits, a queued buyer weighted by fills 2.12\%, 2.64\% and 3.20\%, and the price gives back 1.76\%, 2.81\% and 5.27\% over the following week.
\item At 10\%, the backlog alone gives 3.42\% (1.21\% filled); attention adds most of the rest and takes it back over the week.
\item It needs a short sale, and borrow may be scarce or banned.
\item Quarterly holdings five days late leave an IC of 0.003 and a Sharpe ratio of 0.26, against 0.021 and 2.04 with daily data.
\item Fills from the queue, the open as the exit, stamp duty and commissions, the limits in force at each date, and only borrowable stocks for shorts.
\item The northbound-flow signal, whose data were cut by the exchanges.
\item Both trade the \omterm{def:s1:flows:fit}{price pressure} of predictable demand; here the demand is attention buying and queued orders instead of funds' flows.
\item It sells time: it delays a move to the next day, and those who can queue early or sell into the next morning's buyers collect for it.
\end{enumerate}
\end{solution}

\begin{solution}{iq:s1:asia-specific-equity-strategies:1}
Closes at the limit are not the prices at which the stock would have traded: large moves are split over several days, so closes show continuation
that nobody could buy at those prices. A momentum signal fitted on them finds a return it cannot trade.
\end{solution}

\begin{solution}{iq:s1:asia-specific-equity-strategies:2}
Decide whether the backlog and news justify the queue: joining late means filling only if the lock breaks, which is when the trade is worst. Either
queue early with a size the fills can bear, or wait for the next open and pay the gap.
\end{solution}

\begin{solution}{iq:s1:asia-specific-equity-strategies:3}
Compare the speed of price moves (or the hazard of reaching the next tick) as a function of the distance to the limit with that of the same stocks far
from it, controlling for volatility; the magnet shows as acceleration near the limit, stronger at the upper one.
\end{solution}

\begin{solution}{iq:s1:asia-specific-equity-strategies:4}
Fills at the limit only from the queue, with priority; no trades beyond the limits; no same-day sale of shares bought that day; limits by board and
status as they were at each date; stamp duty on sales.
\end{solution}

\begin{solution}{iq:s1:asia-specific-equity-strategies:5}
A stock locked limit down cannot be sold for as many days as the fall needs; value-at-risk from daily returns understates it; a fund with daily
redemptions may have to sell what it can, not what it should.
\end{solution}

\begin{solution}{iq:s1:asia-specific-equity-strategies:6}
With $y = x - L$ exponential with mean $m$, $E[e^{-y/s}] = s/(s + m)$ and $E[y e^{-y/s}] = m s^2/(s + m)^2$, so the gain of a filled order is
$E[y e^{-y/s}]/E[e^{-y/s}] = m s/(m + s)$, below both $m$ and $s$: fills select the small backlogs, and the unconditional mean $m$ overstates the gain.
\end{solution}
